% SOURCE_AUTHORITY: sga5_fr_workpass.tex, lines 14705--14737 (LF-normalized unit).
% SOURCE_SHA256: 791F4EFFC5E02832D5D77ED03518C8156D6F07E4C8238B03545DB93D883FBB28
Estabeleçamos primeiro a); deduziremos b) e c). Seja $\varphi_F$ o endomorfismo considerado de $H^0(X,F)$; é claro que $\varphi_F$ depende funtorialmente de $F$ e que $\varphi_X=\id_{H^0(X,X)}$. Ora, todo elemento $s$ de $H^0(X,F)$ identifica-se com um morfismo de feixes $s:X\to F$ tal que $s$ é a imagem, por $H^0(X,s)$, do elemento único de $H^0(X,X)$; como
\[
\varphi_F\circ H^0(X,s)=H^0(X,s)\circ\varphi_X=H^0(X,s),
\]
pelas observações precedentes, vê-se que $\varphi_F(s)=s$ para todo $s$, logo $\varphi_F=\id_{H^0(X,F)}$.

Observe que $\varphi_F$ é também o composto
\[
H^0(X,F)\longrightarrow H^0(X,\fr_{X*}F)
\longrightarrow H^0(X,F),
\]
onde a primeira flecha é $H^0(X,F(\Fr_{/X})^{-1})$ e a segunda é a aplicação canônica, que não é senão a identidade. Do mesmo modo, os endomorfismos considerados em b) e c) fatoram-se assim:
\[
H^1(X,F)\longrightarrow H^1(X,\fr_{X*}F)
\longrightarrow H^1(X,F),
\]
e
\[
\RGamma_X(F)\longrightarrow \RGamma_X(\fr_{X*}F)
\longrightarrow \RGamma_X(F).
\]

Para estabelecer b), utilizaremos um monomorfismo $F\to G$ de $F$ em um feixe de grupos $G$ tal que $H^1(X,G)=0$; então $\fr_{X*}F\to\fr_{X*}G$ é um monomorfismo e $H^1(X,\fr_{X*}G)=0$. Introduzamos o feixe quociente $C=G/F$ de espaços homogêneos. Como $\fr_X$ é integral, sobrejetivo e radicial, os funtores $\fr_X^*$ e $\fr_{X*}$ são equivalências de categorias quase-inversas uma da outra (SGAA VIII teorema 1.1); deduz-se que
\[
\fr_{X*}(C)\simeq\fr_{X*}(G)/\fr_{X*}(F),
\]
e a sequência exata de cohomologia (SGAA XII proposição 3.1) mostra que o endomorfismo considerado de $H^1(X,F)$ provém, por passagem ao quociente, de
\[
\varphi_C:H^0(X,C)\longrightarrow H^0(X,C),
\]
que é a identidade de acordo com a); este endomorfismo é, portanto, $\id_{H^1(X,F)}$.

A afirmação c) demonstra-se de maneira análoga, utilizando um complexo $I$ injetivo em cada grau e isomorfo a $F$ na categoria derivada $D^+$, construída sobre os $\Lambda$-módulos; o complexo $\fr_{X*}(I)$ é injetivo em cada grau e isomorfo a $\fr_{X*}(F)$ em $D^+$, pois $\fr_{X*}$ é uma equivalência de categorias. Pode-se então identificar $\RGamma_X(F)$ com $H^0(X,I)$ e $\RGamma_X(\fr_{X*}F)$ com $H^0(X,\fr_{X*}I)$; vê-se que o endomorfismo considerado de $\RGamma_X(F)$ identifica-se com $\varphi_I$, que é a identidade; este endomorfismo é, portanto, $\id_{\RGamma_X(F)}$.
